\section*{The results and their mechanisms}
We write $F_b=f^{-1}(b)$ for the full Cartier fibre and
$S_b=(F_b)_{\mathrm{red}}$ for its reduced support.  All line bundles
and all restriction maps below are holomorphic.  The rational Betti
conditions are $b_1(Y)=b_2(Y)=0$; the \emph{integral CDP hypotheses}
mean $H^1(Y,\Z)=H^2(Y,\Z)=0$, together with the holomorphic
connected-fibre algebraic reduction used in the cited statement.
We keep these two scopes explicit.

\subsection*{1. The omitted cotangent term}
\paragraph{The step as printed.}
On the published p.~692 of \cite{CDP2020}, the authors put
$S=\operatorname{red}X_c$ and display
\[
 0\longrightarrow N_S^*\longrightarrow\Omega_X^1|_S
   \longrightarrow\Omega_S^1\longrightarrow0.
\]
The next sentence is:
\begin{quote}
``it is immediate, provided $\mathcal L|_S$ is non-torsion,
that it suffices to show the statement''
\[
 H^0(S,\widetilde\Omega_S^1\otimes\mathcal L|_S)=0.
 \tag{11}
\]
\end{quote}
Here the printed definition is
$\widetilde\Omega_S^1=\Omega_S^1/\mathrm{torsion}$.
We read this as the asserted sufficiency of the displayed
torsion-free vanishing for the ambient cotangent vanishing.
The following exact-sequence calculation identifies the additional
map whose kernel must vanish.

\RungZero{For a reduced hypersurface $S\subset Y$, torsion in
$\Omega_S^1$ contributes ambient cotangent sections through a
connecting-map kernel.  The two endpoint vanishings on p.~692 of
\cite{CDP2020} leave exactly this kernel; they determine the desired
vanishing only after its injectivity has also been proved
(\ref{cdp:st:A.1.1}).}
\RungOne
\begin{enumerate}[label=\arabic*.,leftmargin=*]
\item Put $N^*=\OO_Y(-S)|_S$, $V=\Omega_Y^1|_S$,
$T=\operatorname{Tors}\Omega_S^1$, and $Q=\Omega_S^1/T$.
The conormal map $d:N^*\to V$ is injective because it is
nonzero at the generic point of every component of the reduced
hypersurface.
\item The inverse image $K$ of $T$ in $V$ gives the exact sequences
$0\to N^*\to K\to T\to0$ and $0\to K\to V\to Q\to0$.
These are the actual maps induced by the conormal sequence.
\item Tensor by a line bundle $A$.  The equality
$H^0(Q\otimes A)=0$ identifies $H^0(V\otimes A)$ with
$H^0(K\otimes A)$ through the second sequence.
\item The equality $H^0(N^*\otimes A)=0$ identifies this last
space with the kernel of
$\delta_A:H^0(T\otimes A)\to H^1(N^*\otimes A)$ through the
first sequence.  This proves the exact missing condition.
\end{enumerate}

\subsection*{2. An explicit nodal counterexample}
\RungZero{A nodal cubic times $\PP^1$ has a one-dimensional
ambient cotangent space for every gluing parameter, while both
proposed endpoint spaces vanish for every nontrivial gluing
parameter.  Normalization preserves the conductor-supported
section responsible for the difference (\ref{cdp:st:B1}).}
\RungOne
\begin{enumerate}[label=\arabic*.,leftmargin=*]
\item The family
$y^2z=x^2(x+z)+tz^3$ over $|t|<\varepsilon<4/27$ has smooth
total space and only one singular fibre in this disc: its central
nodal cubic $C_0$.  The derivative in $t$ and the quadratic term
$y^2-x^2$ verify these assertions.
\item Normalization is $u\mapsto(u^2-1,u(u^2-1))$; its points
$1,-1$ are identified.  The gluing line $A_\lambda$ has descent
complex $\C\xrightarrow{1-\lambda}\C$, so
$h^0(A_\lambda)=h^1(A_\lambda)$ equals one at $\lambda=1$ and
zero elsewhere.
\item At $vw=0$, torsion is generated by $v\,dw=-w\,dv$.
The torsion-free quotient is $\nu_*\OO_{\PP^1}(-2)$ and its
inverse image in the ambient cotangent is $\nu_*\OO_{\PP^1}$.
The last sheaf has one section after every gluing twist.
\item The product decompositions of the cotangent sheaves and
K\"unneth give every entry of the cohomology table in
\ref{cdp:st:B1}; in particular the ambient $h^0$ is always one.
\item The divisor $P-Q$ with normalized coordinates
$u(P)=(1-\lambda)/(1+\lambda)$ and $u(Q)=0$ has gluing scalar
$\lambda$.  Extending the smooth points to sections extends this
line bundle to the given smooth total space.
\end{enumerate}

\subsection*{3. The torsion-free assertion on every reduced support}
\RungZero{The conclusion of \cite[Proposition~7.3, p.~14 of the attached-source edition]{CDP2020}
holds through unitary local systems: the transgression character
has zero first cohomology on every fibre support, and a
torsion-free holomorphic one-form injects into that cohomology
(\ref{cdp:st:A.4.7}--\ref{cdp:st:A.4.9}).}
\RungOne
\begin{enumerate}[label=\arabic*.,leftmargin=*]
\item Integral Leray gives a nonzero transgression section
$\gamma$ of $R^1f_*\Z$.  A meridian bounds a normal disc with
positive multiplicity, which proves that its specialization is
nonzero on every full fibre (\ref{cdp:st:D.1.8}).
\item Coherent Leray lifts its coefficient image to
$\beta\in H^1(Y,\OO_Y)$.  For real irrational $r$,
$L_r=\exp(2\pi i r\beta)$ restricts on every tube to the unitary
character $\exp(2\pi i r\gamma)$, of infinite order.
\item On a very general smooth fibre, the corresponding degree-zero
and degree-one local-system groups vanish for every irrational nonzero
integral character. The torus calculation is the contracted Koszul
complex; finite covers and the explicit group extensions give the
remaining calculations. Proper smooth transport along a path gives a
diffeomorphism to any smooth fibre and identifies its character local
system by pullback, proving the same vanishing there.
\item The punctured-tube Leray sequence gives degree-one
vanishing.  The surjection of its fundamental group onto that of
the full tube makes pullback on degree-one local-system cohomology
injective, and proper base change gives the vanishing on $S_b$.
\item Pull a twisted holomorphic one-form to the normalization
resolution.  Stokes' theorem makes it closed; continuous local
primitives, holomorphic on normalization, define its class in
$H^1(S_b,L_r^{\mathrm{flat}})$.  A zero class has a global
parallel primitive, proving the required injection.
\item Proper semicontinuity for the finitely many singular
supports turns this one simultaneous family of vanishing twists
into a dense open set in $\Pic(Y)$.
\end{enumerate}

\subsection*{4. The component count with monodromy}
\RungZero{For $s$ deleted values on a base of genus $g$, the
number $r$ of irreducible deleted fibre components is
$r=2g+s-1+t'-b_1(Y)$, where $t'$ is the dimension of the
monodromy coinvariants of $H_1(F,\Q)$; the complement records
those coinvariants in its first homology.  The rational formula
is credited to \cite[Section~10.5, p.~91]{AC2026}; the integral
refinement is \ref{cdp:st:D.1.2}.}
\RungOne
\begin{enumerate}[label=\arabic*.,leftmargin=*]
\item The punctured base has free fundamental group of rank
$2g+s-1$.  Lifts of its free generators split the bundle's
fundamental-group extension.
\item Abelianizing gives
$H_1(U,\Z)=\Z^{2g+s-1}\oplus H_1(F,\Z)_{\pi_1(B^\circ)}$.
The dual of the coinvariant quotient is the invariant subspace
of $H^1(F,\Q)$.
\item The compact-support sequence and Poincar\'e duality give
$0\to\Q^r\to H_1(U,\Q)\to H_1(Y,\Q)\to0$ under
$b_2(Y)=0$.  Taking dimensions proves the formula.
\item Under the integral CDP hypotheses, universal coefficients
give $H_1(Y,\Z)=0$.  The integral compact-support sequence gives
$H_1(U,\Z)=\Z^r$; the displayed coinvariant summand is therefore
free.  The injection $H^1(B,\Q)\to H^1(Y,\Q)$ also gives
$B=\PP^1$.
\end{enumerate}

\subsection*{5. The smooth general fibre}
\RungZero{Under the integral CDP hypotheses the very general
fibre is a two-torus or a Hopf surface of algebraic dimension
zero.  Adjunction and divisor Riemann--Roch restrict the surface
invariants; the corrected integral component count and the
explicit monodromy comparisons remove the remaining entries
(\ref{cdp:st:C.3.1}).}
\RungOne
\begin{enumerate}[label=\arabic*.,leftmargin=*]
\item The normal line of a regular fibre is trivial.
Restriction of ambient Chern classes and the divisor sequence give
$c_1(F)=0$ rationally and $\chi(\OO_F)=0$.
\item The surface classification at these invariants leaves the
explicit list stated in \ref{cdp:st:C.3.1}; its listed surfaces
have Euler number zero by their Chern numbers.
\item The integral coinvariant quotient in
\ref{cdp:st:D.1.2}, applied to the explicit finite and nilpotent
surface groups, gives the exclusions proved in
\ref{cdp:st:C.3.7}, \ref{cdp:st:C.3.9}, and
\ref{cdp:st:C.3.10}.  The Inoue exclusion is the further
monodromy comparison in \ref{cdp:st:C.3.12}.
\item An elliptic Hopf very general fibre supplies a moving
relative anticanonical system.  Its relative algebraic reduction
would raise the algebraic dimension of $Y$ above one; the
construction of this system in \ref{cdp:st:C.3.8} proves the
contradiction.
\item The quotient constructions of \ref{cdp:st:C.3.6} realize
both primary and secondary Hopf smooth fibres, with the stated
topology, so these surviving possibilities have actual examples.
\end{enumerate}

\subsection*{6. The sign for irreducible and normal-crossings supports}
\RungZero{For a torus fibration with $b_1(Y)=b_2(Y)=0$ whose
reduced supports are irreducible or have normal crossings, each
support has nonnegative Euler number.  Summing these numbers and
using $e(Y)=2-b_3(Y)$ gives $c_3(Y)\in\{0,2\}$; the rational
normalization contributes two and the other irreducible cases
contribute zero (\ref{night:P4.58}).}
\RungOne
\begin{enumerate}[label=\arabic*.,leftmargin=*]
\item Connected-fibre Leray gives $B=\PP^1$.  The global
transgression restricts to a nonzero integral class on every
support, which supplies the finite-character comparisons used in
its normalization calculation.
\item For a normal-crossings support, adjunction on each
normalized double curve gives $e(C_\alpha)=\deg R_\alpha$,
with three marked branches at every triple point.  The weighted
nearby-fibre Euler formula then gives $e(S)=\tau(S)$
(\ref{cdp:st:C.2.NC.2}--\ref{cdp:st:C.2.NC.4}).
\item For an irreducible singular support with rational
normalization resolution, the finite-character support comparison
gives $e(S)=2$ (\ref{night:P4.35}).
\item A class-VII normalization resolution reduces to a smooth
minimal Hopf normalization through the effective conductor and
curve-image comparisons of \ref{night:P4.36}--\ref{night:P4.37}.
Its Euler number is zero.
\item In the ruled case, the remaining positive-Euler conductor
configuration has a doubled horizontal pole section.  The local
intersection-form comparison excludes its contraction, and the
spectral-cover and conductor-graph calculations exclude each
noncontracted degree-two and degree-three alternative
(\ref{night:P4.54}--\ref{night:P4.57}, \ref{night:P4R.20}).
The surviving normalization is smooth minimal elliptic ruled,
with Euler number zero.
\item Euler additivity gives $c_3(Y)=\sum_b e(S_b)\ge0$.
Poincar\'e duality and the skew middle pairing give
$e(Y)=2-b_3(Y)$ with $b_3(Y)$ even, proving the two possible
values and the equality criterion.  Paper~VI treats arbitrary
reducible supports with their multiplicities.
\end{enumerate}

\subsection*{7. The exact replacement for the generic vanishing claim}
\RungZero{Under the integral CDP hypotheses, the generic
vanishing $H^q(Y,T_Y\otimes L)=0$ for $q\ne1$ holds exactly
when the Picard pullback to the normalization of every irreducible
nonnormal critical support is nonzero.  A zero pullback preserves
the conductor section for every twist and forces
$H^2(Y,T_Y\otimes L)\ne0$ for every $L$
(\ref{cdp:st:C4f}).}
\RungOne
\begin{enumerate}[label=\arabic*.,leftmargin=*]
\item For every critical full fibre, its ambient-section
nonvanishing locus is a closed analytic subset of $\Pic(Y)$.
The finite full-fibre test of \ref{cdp:st:C4c} and its
transgression version \ref{cdp:st:C4d} characterize generic
degree-two vanishing by properness of these loci.
\item Smooth supports have proper loci by the cyclic-cover and
nilpotent-layer calculation.  Reducible supports have proper loci
by successively twisting against an adjacent component and then
using the full Cartier filtration (\ref{cdp:st:A.3.5}).
\item Persistent nonvanishing on an irreducible support gives
sections of one fixed line on a normalization resolution after
every ambient twist.  Gauduchon degree forces those pullbacks to
be unitary-flat (\ref{cdp:st:A.5.4}).
\item The conductor relation and the exact curve-genus bounds
then make the pullback homomorphism zero.  In the final ruled
case the inverse-Jacobian line has sections at at most two points
of the elliptic Picard group; a connected persistent image is
therefore the identity (\ref{cdp:st:C4e}).
\item Conversely, a zero Picard pullback supplies a nonzero
normalization section in the last ideal layer of $mS$.
The explicit injection \eqref{cdp:eq:C.VC.4} puts it in the
ambient cotangent on the full fibre, for every twist.
\item Fibre duality, Leray, and semicontinuity give the generic
criterion and the all-twist nonvanishing.  These maps have the
scope stated in \ref{cdp:st:C4f}; their individual ranks in the
standard family are computed in the next result.
\end{enumerate}

\subsection*{8. Integral twists and the rank of the normalization map}
\RungZero{The tangent cohomology of every integral twist of
$X$ follows from three split root-curve direct images and one
length-two root-chart module.  The map to
$H^2(\mathrm dP_6,K)=\C$ has rank one exactly for
$n\in\{-6,-3,-2\}\cup\Z_{\ge0}$; at $n=-1$ its source is
one-dimensional and the map is zero (\ref{cdp:st:C5b}).}
\RungOne
\begin{enumerate}[label=\arabic*.,leftmargin=*]
\item On $\mathfrak B=\PP(3,4)$, the vertical Lie and
structure-sheaf bundles are $E=G=\OO\oplus\OO(-1)$.
The logarithmic tangent sequence has quotient $\bar f^*\OO(-5)$.
\item Its two cup-product boundary matrices are
$\delta_0=(0,y,y,x)^t$ and
$\delta_1=\left(\begin{smallmatrix}-y&0\\-x&y\end{smallmatrix}\right)$,
with $x,y$ of weights $3,4$.  Their exact kernels and cokernels
give the direct images in \eqref{cdp:eq:C.8}--\eqref{cdp:eq:C.9}.
\item Removing the finite stabilizer characters yields
$\rho_*\OO(n)=\OO(d(n))$, with $d(n)$ in
\ref{cdp:st:C5b}.  Invariants in the length-two module have
dimension $0$ for $3\mid n$ and $1$ otherwise.
\item Cohomology of these line bundles on $\PP^1$ and its
two-column Leray sequence give the four exact dimensions
\eqref{cdp:eq:C.6}.
\item The normalization quotient on the root curve is
$0\to\OO(n-6)\to\OO(n+6)\to\C_{p_0}\to0$.
Evaluation at $p_0$ is nonzero exactly when
$n+6\in\langle3,4\rangle$, giving the stated rank set.
\end{enumerate}

\subsection*{9. Trivial monodromy and an equality example}
\RungZero{Trivial rational monodromy forces $c_3(Y)\le0$ in
the torus-fibred Betti class, because a persistent positive-index
obstruction would resolve into components of geometric genus
zero and force nontrivial nearby monodromy.  A holomorphic
three-piece filling realizes equality with
$H_1(Y,\Z)=H_2(Y,\Z)=0$
(\ref{cdp:st:D.2.5}, \ref{cdp:st:D.2.7}).}
\RungOne
\begin{enumerate}[label=\arabic*.,leftmargin=*]
\item Positive Euler number makes $b_3=0$, and the index and
finite-fibre test give one full fibre carrying ambient sections
for every Picard twist.
\item The Cartier filtration selects one persistent graded
piece.  Its pullback to the normalization resolution has degree
zero after every ambient twist.  The effective conductor then
forces $p_g=0$ on that resolution.
\item Embedded resolution adds only ruled surfaces and projective
planes.  The local period comparison of \ref{cdp:st:D.2.3}
for a torus degeneration with these $p_g=0$ components yields
nontrivial $H^1$ monodromy, proving the inequality.
\item For equality, set $E=\C/(\Z+i\Z)$ and
$U=\OO_{\operatorname{Bl}_{(0,0)}(E\times\Delta)}(R)^\times/\langle e^{-2\pi}\rangle$.
This proper local model has central fibre a primary Kodaira
surface and a primary Hopf surface meeting in $E$; its nearby
fibre is $E\times E$ with the specified product marking.
\item Three such pieces glued by the translations in
\eqref{cdp:eq:D.28} give the exact Mayer--Vietoris quotients
\eqref{cdp:eq:D.31}--\eqref{cdp:eq:D.32}.  Their displayed
integral relation bases give $H_1=H_2=0$.
\item Every central component and double curve has Euler number
zero, so $e(Y)=0$ and $b_3(Y)=2$.  Restricting an effective
divisor to a general torus excludes horizontal divisors and
proves $\mathcal M(Y)=f^*\C(t)$, hence $a(Y)=1$.
\end{enumerate}

\subsection*{10. Retention of the fibration under deformation}
\RungZero{Rank two of the torus Kodaira--Spencer tensor at one
regular value makes every small deformation retain the
holomorphic fibration.  The same cup-product map kills each
Artin-order map obstruction, and the relative graph space turns
these lifts into a holomorphic family, including nonreduced bases
(\ref{cdp:st:E6}).}
\RungOne
\begin{enumerate}[label=\arabic*.,leftmargin=*]
\item The tangent sequence on a rank-two fibre makes
$H^1(T_F)\to H^1(T_Y|_F)$ surjective: its next boundary is
the invertible cup-product map of \ref{cdp:st:E5}.
\item A torus has a smooth local deformation space.  Across a
small Artin extension, changing a lifted torus by a class in
$H^1(T_F)$ kills its embedding obstruction in $H^1(T_Y|_F)$.
\item The obstruction to lifting the global map is a class in
$H^1(Y,f^*T_{\PP^1})$.  The lifted Cartier fibre makes its
restriction zero at every rank-two regular value.
\item Leray identifies that class with a section of
$T_{\PP^1}\otimes R^1f_*\OO_Y$.  Local freeness and vanishing
on an open set make the section zero, so the map lifts at every
Artin order.
\item The relative Douady graph germ represents these maps.
The lifting property makes it smooth over the deformation base;
a local section gives the required analytic map.
\item Connectedness and torus type persist near a regular fibre.
Stein factorization then gives connected fibres globally; under
$b_2(Y)=0$, the vertical-divisor argument gives $a(Y_s)=1$.
The invertible logarithmic cusp coefficient verifies the rank
condition for every named family in \ref{cdp:st:E7}.
\end{enumerate}

\subsection*{11. The two-parameter, Liu--Rollenske and positive-comb germs}
\RungZero{The two-parameter and Liu--Rollenske constructions
have tangent cohomology $(1,2,2,0)$ and smooth two-dimensional
Kuranishi germs.  A positive $m$-comb has
$(2m+1,8m+1,6m+1,0)$ and a smooth germ of dimension $8m+1$:
root-curve cohomology gives the dimensions, and actual period
sections supply all tangent directions
(\ref{cdp:st:E13}--\ref{cdp:st:E15}).}
\RungOne
\begin{enumerate}[label=\arabic*.,leftmargin=*]
\item Retain the finite characters on the root curve.  The
explicit logarithmic connection and normalization cup products
identify the vertical Lie, Hodge-quotient, and determinant bundles.
\item Compute their cohomology by invariant generators at each
root point.  The missing product generators give the precise
skyscraper corrections in \eqref{cdp:eq:E15-products}.
\item The rank-two cusp tensor makes the relevant cup-product
map injective.  The normal quotient at the Mumford fibres and its
residue map give the displayed tangent tuples in the full
logarithmic tangent sequences.
\item For the two-parameter cases, the exceptional-value
configuration and the double-curve period detect the two
independent Kodaira--Spencer directions.
\item For a positive comb, the normal bundle of the period
section has $H^1=0$ and $h^0=8m+1$.  The finite root conditions
and the actual cusp gluings therefore give a smooth family of
that dimension.
\item A trivial total-space deformation induces a base vector
field vanishing at all three distinguished values and then a
zero fibrewise period derivative.  This proves tangent
injectivity; the resulting tangent isomorphism makes the
Kuranishi defining ideal zero.
\end{enumerate}

\subsection*{12. Separation along analytic and continuous deformations}
\RungZero{The ordered coprime pair $(3k,4l)$ is constant in every
proper smooth holomorphic family over a connected reduced analytic
base containing a constructed member, and it distinguishes the
constructed members' deformation classes through arbitrary finite
chains.  Disc propagation and finite curve projections make the
fixed-pair locus open and closed; the actual Kuranishi retraction
gives the same separation in the $C^\infty$ space of integrable
complex structures and its indicated diffeomorphism quotients
(\ref{push:src:E.2.23}--\ref{push:src:E.2.24}).}
\RungOne
\begin{enumerate}[label=\arabic*.,leftmargin=*]
\item The fibration-retention theorem supplies the pencil near
each constructed member; its intrinsic anticanonical description
identifies the pencils on overlapping deformation charts.
\item The finite covers carry the labelled integral monodromies
and attachment classes.  Their holomorphic continuation preserves
the specified root charts and finite multiplicities.
\item The marked period map has the given root-curve form on the
known open set.  The analytic comparisons in
\ref{push:src:E.2.21} extend a full punctured-disc family, and
\ref{push:src:E.2.22} fills its successive discrete exceptional
sets in reverse order to propagate from any nonempty known open
subset of a disc.
\item The cusp comparison retains the nonzero logarithmic
coefficient and the full affine gluing class.  Together with the
finite charts, it identifies the continued members with the same
$(3k,4l)$ construction.
\item The fixed-pair locus is open by local miniversality.  At a
boundary point $b$, choose a finite local projection
$\rho:T\to\Delta^d$ of an irreducible branch with
$\rho^{-1}(0)=\{b\}$.  The inverse image of the line through
$0$ and $\rho(a)$, for a known point $a$, has an irreducible
curve containing both $a$ and $b$; pullback to its normalization
and the disc propagation just proved show that $b$ is known
(\ref{push:src:E.2.23}).
\item The locus is therefore closed, and connectedness makes it
the whole analytic base.  This preserves the pair through every
finite deformation chain; conversely the admissible half-plane
and its quotient by the actual period lattice $2kl\Z$ give a
connected family for each fixed pair
(\ref{push:src:E.2.23}).
\item At every smooth integrable structure $J_0$, the Kuranishi
local theorem supplies a connected neighbourhood $\mathcal U$
and a retraction $\Xi:\mathcal U\to K$, with each $J$
biholomorphic to $\Xi(J)$ through a diffeomorphism isotopic to
the identity.  Apply the analytic-base conclusion to the
connected reduction of $K$: if $\mathcal U$ meets the fixed-pair
locus, every point of $\mathcal U$ belongs to it.  Thus that locus
is open and closed, as are its images under the specified
saturated diffeomorphism quotients, and every continuous family
over a connected parameter space retains its pair
(\ref{push:src:E.2.24}).
\end{enumerate}

\subsection*{13. Four explicit corrections}
\RungZero{The Poincar\'e line bundle gives a torsion higher
direct image in the scope of the printed local-freeness claim;
a square-zero embedded point survives every divisor twist;
positive combs have $2m+1$ global vector fields; and an explicit
half-period section makes the proposed theta square vanish
while the actual period determinant is nonzero.  The four
constructions and their exact corrections are
\ref{cdp:st:E12}, \ref{cdp:st:E17}, \ref{inv:E2.4}, and
\ref{cdp:st:E18}, with their source scopes retained.}
\RungOne
\begin{enumerate}[label=\arabic*.,leftmargin=*]
\item For the normalized Poincar\'e line $P$ on $E\times E$,
the moving degree-one divisor gives
$Rp_*P=\C_0[-1]$, with length one at the origin.
Tensoring with $R\Gamma(\PP^1,\OO(-2))=\C[-1]$ gives the
direct-image counterexample $\C_0[-2]$.
\item On $\OO_Z=\OO_{\PP^1}\oplus\C_p$ with square-zero
second summand and a Cartier divisor $q\ne p$, the equality
$\OO_Z(-kq)=\OO_{\PP^1}(-kq)\oplus\C_p$ retains a nonzero
section for every $k$.
\item Krull intersection proves the corrected stable-space
formula: descending twists stabilize at the sections whose
support avoids the Cartier divisor.  Absence of embedded
components gives the stated eventual vanishing case.
\item The computed $h^0(T_{X_m})=2m+1$ is at least $25$ for
$m\ge12$.  Its comparison with the printed two-vector-field
bound gives the explicit counterexample and the exact
algebraic-dimension scope described in \ref{inv:E2.4}.
\item On a small disc about $\tau=i$, take $\mu=1/2$ and
\[
 v=\frac{i}{\pi}\log\frac{\vartheta_1(1/2,\tau)}{\eta(\tau)}+i,
 \qquad
 \Lambda_\tau=\Z(1,0)+\Z(0,1)+\Z(\tau,1/2)+\Z(3,6v).
\]
The logarithm is real at $i$; the theta products give
$\operatorname{Im}v>0$, so $\C^2/\Lambda_\tau$ is a proper
family of compact tori.  The quotient
$\vartheta_1(z-1/2,\tau)/\vartheta_1(z,\tau)$ has divisor
$[1/2]-[0]$ and multipliers $1,-1$, giving the actual elliptic
line bundle and its contracting linearization.  Adjoining
$(1/2,v)$ enlarges the lattice by index six; identity on
$\C^2$ induces the degree-six isogeny to the symmetric-period
family with $\Omega=\left(\begin{smallmatrix}\tau&1/2\\1/2&v\end{smallmatrix}\right)$
and preserves its Kodaira--Spencer rank
(\ref{cdp:st:E18}).
\item Put $h=\partial_\mu\log\vartheta_1(\mu,\tau)$.
The half-period shift and evenness of $\vartheta_2$ give
$h=\mu'=0$, whereas the heat equation gives
\[
 \det\Omega'(i)=-\frac{\wp(1/2;i)}{4\pi^2}
 =-\frac{\vartheta_3(0,i)^4+\vartheta_4(0,i)^4}{12}\ne0.
\]
In the displayed linearization the full identity is
\[
 \det\Omega'
 =-\left(\mu'-\frac{i}{2\pi}h\right)^2
       -\frac{\wp(\mu,\tau)}{4\pi^2}
 \qquad\text{(\ref{cdp:eq:E-theta-general}).}
\]
For the actual original period branch, the equation
$\wp(\mu,\tau)=0$ proves its square formula; an added
holomorphic linearization term contributes its derivative
(\ref{cdp:st:E18}).
\end{enumerate}

\subsection*{14. The base parameter of the Liu--Rollenske family}
\RungZero{For the fixed attachment data, the family closes
after the parameter is lifted through $u=w^2$.  The odd return
changes the multiple fibre's Albanese curve, so it does not
identify the two $w$-values; the tangent-cohomology calculation
remains the one stated in \ref{cdp:st:E14}
(\ref{ii:LR:covering}).}
\RungOne
\begin{enumerate}[label=\arabic*.,leftmargin=*]
\item The explicit finite-fibre lattice in
\ref{ii:LR:covering} gives its Albanese elliptic curve with the
specified period parameter $q$.
\item The return of the scalar parameter sends $q$ to $-q$.
The resulting lattices and their modular invariants are compared
by the morphisms in that statement.
\item The odd return therefore changes the biholomorphism
invariant of the member with the fixed data.  The even return
preserves the full lattice and attachment maps, yielding the
family over the punctured $w$-disc covering the punctured $u$-disc.
\item The logarithmic tangent and residue calculations use these
fixed local models, so the proved tuple $(1,2,2,0)$ and its
local Kuranishi conclusion are retained.
\end{enumerate}

\subsection*{15. The remaining comb splitting data}
\RungZero{The tangent cohomology and Kuranishi germ of a
positive comb are determined above.  The individual splitting
of the rank-three kernel $K_m$ still requires the actual
period-section map: the supplied existence family does not
specify the transition matrices of one smoothed $12$-tooth
member (\ref{inv:E2.3}).}
\RungOne
\begin{enumerate}[label=\arabic*.,leftmargin=*]
\item The defining kernel sequence gives
$\deg K_m=-2m-2+e_m$, where $e_m$ is the cokernel length of
the specified map; ranks and degrees prove this identity.
\item The support and elementary divisors of that cokernel
depend on the actual period section, so the degree identity
retains those finite data as inputs.
\item The period-section family and its tangent isomorphism
prove the Kuranishi statement independently of this individual
splitting.  The final part records the exact unsolved splitting
instance and every established constraint on it.
\end{enumerate}
